% Algebra 2 Lesson 2-2 fast-track packet: Standard Form of a Quadratic Function
% (enVision A2 2024 TE, Example 2 and Example 3 only; draft 2026-10-30, PLC has not met on this lesson)
% Shared preamble for the "fast" (hacked-down) Algebra 2 lesson packets, Oct 2026.
% Each packet: % Shared preamble for the "fast" (hacked-down) Algebra 2 lesson packets, Oct 2026.
% Each packet: % Shared preamble for the "fast" (hacked-down) Algebra 2 lesson packets, Oct 2026.
% Each packet: \input{fast-packet-template} then \begin{document} ... \end{document}.
\documentclass[12pt]{article}
\usepackage[margin=0.7in]{geometry}
\usepackage{amsmath,amssymb}
\usepackage{booktabs,tabularx,array}
\usepackage{enumitem}
\usepackage[T1]{fontenc}
\usepackage[utf8]{inputenc}
\usepackage{lmodern}
\usepackage{tikz}
\usepackage{pgfplots}
\pgfplotsset{compat=1.18}
\usepackage{xcolor}
\setlist{nosep,leftmargin=1.4em}
\setlength{\parindent}{0pt}
\setlength{\parskip}{6pt}
\renewcommand{\arraystretch}{1.35}
\newcolumntype{Y}{>{\raggedright\arraybackslash}X}
\pagestyle{empty}

% blanks and workspace
\newcommand{\blank}[1][2.6cm]{\rule[-2pt]{#1}{0.5pt}}
\newcommand{\lines}[1][1]{\par\vspace{2pt}\foreach \i in {1,...,#1}{\rule{0pt}{22pt}\rule{\linewidth}{0.4pt}\par}}
\newcommand{\wspart}[2]{\vspace{4pt}{\large\bfseries #1: #2}\par\vspace{-6pt}\rule{\textwidth}{0.4pt}\par}

% header: \packethead{Lesson 1-3}{Piecewise-Defined Functions}
\newcommand{\packethead}[2]{%
{\LARGE\bfseries #2}\hfill Name \blank[5cm]\\[2pt]
Algebra 2 \quad #1 (fast track) \hfill Date \blank[2.5cm] \quad Period \blank[1.2cm]\par\vspace{2pt}}

% boxed goals: \goals{learning goal}{language goal}
\newcommand{\goals}[2]{%
\fbox{\parbox{\dimexpr\textwidth-2\fboxsep-2\fboxrule}{%
\textbf{Learning goal:} #1\\[2pt]
\textbf{Language goal:} #2}}\par}

% quiz target box: \quiztarget{text}
\newcommand{\quiztarget}[1]{%
\vspace{4pt}\colorbox{gray!12}{\parbox{\dimexpr\textwidth-2\fboxsep}{\textbf{On the quiz:} #1}}\par}

% IXL table row helper: code & skill & which quiz question it feeds
\newcommand{\ixlrow}[3]{\texttt{#1} & #2 & #3 \\}

% one grid style for every graph; override xmin/xmax/... per plot
\pgfplotsset{sheetgrid/.style={
  axis lines=middle, grid=both, grid style={gray!45},
  xmin=-10.5, xmax=10.5, ymin=-10.5, ymax=10.5,
  xtick={-10,-8,...,10}, ytick={-10,-8,...,10},
  x=0.28cm, y=0.28cm,
  ticklabel style={font=\scriptsize, fill=white, inner sep=1pt},
  xlabel={$x$}, ylabel={$y$},
  axis line style={thick}, clip=true}}
% unit grid for small windows
\pgfplotsset{unitgrid/.style={sheetgrid,
  xmin=-5.5, xmax=5.5, ymin=-5.5, ymax=5.5,
  xtick={-5,...,5}, ytick={-5,...,5}, x=0.5cm, y=0.5cm}}
% open and closed endpoints
\tikzset{closed/.style={circle, fill, inner sep=1.6pt},
         open/.style={circle, draw, fill=white, inner sep=1.4pt}}
 then \begin{document} ... \end{document}.
\documentclass[12pt]{article}
\usepackage[margin=0.7in]{geometry}
\usepackage{amsmath,amssymb}
\usepackage{booktabs,tabularx,array}
\usepackage{enumitem}
\usepackage[T1]{fontenc}
\usepackage[utf8]{inputenc}
\usepackage{lmodern}
\usepackage{tikz}
\usepackage{pgfplots}
\pgfplotsset{compat=1.18}
\usepackage{xcolor}
\setlist{nosep,leftmargin=1.4em}
\setlength{\parindent}{0pt}
\setlength{\parskip}{6pt}
\renewcommand{\arraystretch}{1.35}
\newcolumntype{Y}{>{\raggedright\arraybackslash}X}
\pagestyle{empty}

% blanks and workspace
\newcommand{\blank}[1][2.6cm]{\rule[-2pt]{#1}{0.5pt}}
\newcommand{\lines}[1][1]{\par\vspace{2pt}\foreach \i in {1,...,#1}{\rule{0pt}{22pt}\rule{\linewidth}{0.4pt}\par}}
\newcommand{\wspart}[2]{\vspace{4pt}{\large\bfseries #1: #2}\par\vspace{-6pt}\rule{\textwidth}{0.4pt}\par}

% header: \packethead{Lesson 1-3}{Piecewise-Defined Functions}
\newcommand{\packethead}[2]{%
{\LARGE\bfseries #2}\hfill Name \blank[5cm]\\[2pt]
Algebra 2 \quad #1 (fast track) \hfill Date \blank[2.5cm] \quad Period \blank[1.2cm]\par\vspace{2pt}}

% boxed goals: \goals{learning goal}{language goal}
\newcommand{\goals}[2]{%
\fbox{\parbox{\dimexpr\textwidth-2\fboxsep-2\fboxrule}{%
\textbf{Learning goal:} #1\\[2pt]
\textbf{Language goal:} #2}}\par}

% quiz target box: \quiztarget{text}
\newcommand{\quiztarget}[1]{%
\vspace{4pt}\colorbox{gray!12}{\parbox{\dimexpr\textwidth-2\fboxsep}{\textbf{On the quiz:} #1}}\par}

% IXL table row helper: code & skill & which quiz question it feeds
\newcommand{\ixlrow}[3]{\texttt{#1} & #2 & #3 \\}

% one grid style for every graph; override xmin/xmax/... per plot
\pgfplotsset{sheetgrid/.style={
  axis lines=middle, grid=both, grid style={gray!45},
  xmin=-10.5, xmax=10.5, ymin=-10.5, ymax=10.5,
  xtick={-10,-8,...,10}, ytick={-10,-8,...,10},
  x=0.28cm, y=0.28cm,
  ticklabel style={font=\scriptsize, fill=white, inner sep=1pt},
  xlabel={$x$}, ylabel={$y$},
  axis line style={thick}, clip=true}}
% unit grid for small windows
\pgfplotsset{unitgrid/.style={sheetgrid,
  xmin=-5.5, xmax=5.5, ymin=-5.5, ymax=5.5,
  xtick={-5,...,5}, ytick={-5,...,5}, x=0.5cm, y=0.5cm}}
% open and closed endpoints
\tikzset{closed/.style={circle, fill, inner sep=1.6pt},
         open/.style={circle, draw, fill=white, inner sep=1.4pt}}
 then \begin{document} ... \end{document}.
\documentclass[12pt]{article}
\usepackage[margin=0.7in]{geometry}
\usepackage{amsmath,amssymb}
\usepackage{booktabs,tabularx,array}
\usepackage{enumitem}
\usepackage[T1]{fontenc}
\usepackage[utf8]{inputenc}
\usepackage{lmodern}
\usepackage{tikz}
\usepackage{pgfplots}
\pgfplotsset{compat=1.18}
\usepackage{xcolor}
\setlist{nosep,leftmargin=1.4em}
\setlength{\parindent}{0pt}
\setlength{\parskip}{6pt}
\renewcommand{\arraystretch}{1.35}
\newcolumntype{Y}{>{\raggedright\arraybackslash}X}
\pagestyle{empty}

% blanks and workspace
\newcommand{\blank}[1][2.6cm]{\rule[-2pt]{#1}{0.5pt}}
\newcommand{\lines}[1][1]{\par\vspace{2pt}\foreach \i in {1,...,#1}{\rule{0pt}{22pt}\rule{\linewidth}{0.4pt}\par}}
\newcommand{\wspart}[2]{\vspace{4pt}{\large\bfseries #1: #2}\par\vspace{-6pt}\rule{\textwidth}{0.4pt}\par}

% header: \packethead{Lesson 1-3}{Piecewise-Defined Functions}
\newcommand{\packethead}[2]{%
{\LARGE\bfseries #2}\hfill Name \blank[5cm]\\[2pt]
Algebra 2 \quad #1 (fast track) \hfill Date \blank[2.5cm] \quad Period \blank[1.2cm]\par\vspace{2pt}}

% boxed goals: \goals{learning goal}{language goal}
\newcommand{\goals}[2]{%
\fbox{\parbox{\dimexpr\textwidth-2\fboxsep-2\fboxrule}{%
\textbf{Learning goal:} #1\\[2pt]
\textbf{Language goal:} #2}}\par}

% quiz target box: \quiztarget{text}
\newcommand{\quiztarget}[1]{%
\vspace{4pt}\colorbox{gray!12}{\parbox{\dimexpr\textwidth-2\fboxsep}{\textbf{On the quiz:} #1}}\par}

% IXL table row helper: code & skill & which quiz question it feeds
\newcommand{\ixlrow}[3]{\texttt{#1} & #2 & #3 \\}

% one grid style for every graph; override xmin/xmax/... per plot
\pgfplotsset{sheetgrid/.style={
  axis lines=middle, grid=both, grid style={gray!45},
  xmin=-10.5, xmax=10.5, ymin=-10.5, ymax=10.5,
  xtick={-10,-8,...,10}, ytick={-10,-8,...,10},
  x=0.28cm, y=0.28cm,
  ticklabel style={font=\scriptsize, fill=white, inner sep=1pt},
  xlabel={$x$}, ylabel={$y$},
  axis line style={thick}, clip=true}}
% unit grid for small windows
\pgfplotsset{unitgrid/.style={sheetgrid,
  xmin=-5.5, xmax=5.5, ymin=-5.5, ymax=5.5,
  xtick={-5,...,5}, ytick={-5,...,5}, x=0.5cm, y=0.5cm}}
% open and closed endpoints
\tikzset{closed/.style={circle, fill, inner sep=1.6pt},
         open/.style={circle, draw, fill=white, inner sep=1.4pt}}

\begin{document}

\packethead{Lesson 2-2}{Standard Form of a Quadratic Function}

\goals{I can graph a quadratic function in standard form from its vertex, axis of symmetry and $y$-intercept, and use the vertex to find a maximum or minimum in a real-world problem.}{I can explain how I found a maximum or minimum using the words \emph{standard form}, \emph{vertex}, \emph{axis of symmetry}, \emph{$y$-intercept} and \emph{maximum / minimum}.}

\quiztarget{(1) Find the minimum point of a parabola from its equation (Example 2, Step 1) (full Savvas quiz only). (2) Find the price change that gives a toy maker his highest profit, and that profit (Example 3; bonus). (3) Find the maximum height of a pebble tossed from a cliff (Example 3, Try It 3). (4) Pick the equation of a parabola through three points (full Savvas quiz only). (5) Graph a diver's path from the diving board to the water (Example 2, Try It 2, Practice).}
\topictarget{2}{Question 5 asks you to graph a rock dropped into water (Example 2). Question 6 asks you to graph a launch from a tower and give its maximum height (Examples 2 and 3). Question 3 (bottom of a glass) comes from Example 1 and Question 13 (parabola through the origin and two matching points) from Example 4, which this packet skips.}

\wspart{Example 2}{Graph a Quadratic Function in Standard Form}
How can you use key features to graph $f(x)=x^2-4x+8$? \quad \textbf{Standard form:} $y=ax^2+bx+c$.

\begin{minipage}[t]{0.57\textwidth}
\vspace{0pt}
Identify $a$, $b$ and $c$: \ $a=1$, \ $b=-4$, \ $c=8$.
\begin{enumerate}[label=\textbf{Step \arabic*}, leftmargin=3.4em, itemsep=3pt]
\item \textbf{Vertex and axis of symmetry.}\\
$h=-\dfrac{b}{2a}=-\dfrac{(-4)}{2(1)}=2$\\[2pt]
$k=f(2)=(2)^2-4(2)+8=4$\\
Vertex $(2,4)$; axis of symmetry $x=2$.
\item \textbf{$y$-intercept.} $f(0)=0-0+8=8$, so $(0,8)$. \ (It is always $(0,c)$.)
\item \textbf{Symmetric point.} $(0,8)$ is 2 units \emph{left} of $x=2$, so $(4,8)$, 2 units \emph{right}, is on the parabola too.
\item \textbf{Sketch.} Plot the three points; $a=1>0$, so the parabola opens up and the vertex is the \textbf{minimum}.
\end{enumerate}
\end{minipage}\hfill
\begin{minipage}[t]{0.40\textwidth}
\vspace{0pt}
\begin{tikzpicture}
\begin{axis}[sheetgrid, xmin=-3, xmax=7, ymin=-4, ymax=18,
  xtick={-2,2,4,6}, ytick={4,8,12,16}, minor xtick={-3,...,7}, minor ytick={-4,-2,...,18},
  x=0.6cm, y=0.3cm]
\addplot[very thick, red, domain=-1.1:5.1, samples=80, <->] {x^2-4*x+8};
\node[font=\scriptsize, anchor=north east, inner sep=1.5pt] at (axis cs:0,0) {$O$};
\addplot[thick, red, dashed, <->] coordinates {(2,-3.5) (2,17.5)};
\node[closed, red] at (axis cs:0,8) {};
\node[closed, red] at (axis cs:2,4) {};
\node[closed, red] at (axis cs:4,8) {};
\node[font=\scriptsize, fill=white, inner sep=1pt, anchor=east] at (axis cs:-0.2,8) {$(0,8)$};
\node[font=\scriptsize, fill=white, inner sep=1pt, anchor=west] at (axis cs:4.25,8) {$(4,8)$};
\node[font=\scriptsize, fill=white, inner sep=1pt, anchor=north west] at (axis cs:2.2,3.4) {$(2,4)$};
\node[font=\scriptsize, red, fill=white, inner sep=1pt, anchor=west] at (axis cs:2.2,15.5) {$x=2$};
\end{axis}
\end{tikzpicture}
\end{minipage}

\textbf{Say it:} The \textbf{vertex} is $(\blank[0.8cm],\blank[0.8cm])$ and the \textbf{axis of symmetry} is $x=\blank[0.8cm]$. The \textbf{$y$-intercept} $(0,8)$ and the point $(\blank[0.8cm],\blank[0.8cm])$ are the same distance from the axis of symmetry, so \dots\lines[3]

\newpage

\wspart{Try It 2}{Graph using key features}
Use the key features to graph the function $f(x)=x^2-6x-1$.

\begin{minipage}[t]{0.50\textwidth}
\vspace{0pt}
$a=\blank[1cm]$ \quad $b=\blank[1cm]$ \quad $c=\blank[1cm]$

\vspace{6pt}
$h=-\dfrac{b}{2a}=\blank[3.5cm]$

\vspace{6pt}
$k=f(h)=\blank[4.2cm]$

\vspace{6pt}
Vertex: $(\blank[1cm],\blank[1cm])$

\vspace{6pt}
Axis of symmetry: $x=\blank[1.2cm]$

\vspace{6pt}
$y$-intercept: $(\blank[1cm],\blank[1cm])$

\vspace{6pt}
Symmetric point: $(\blank[1cm],\blank[1cm])$

\vspace{6pt}
Opens \blank[1.6cm], so the vertex is a \blank[2cm].
\end{minipage}\hfill
\begin{minipage}[t]{0.47\textwidth}
\vspace{0pt}
\begin{tikzpicture}
\begin{axis}[sheetgrid, xmin=-2, xmax=8, ymin=-12, ymax=8,
  xtick={2,4,6}, ytick={-8,-4,4}, minor xtick={-2,...,8}, minor ytick={-12,-10,...,8},
  x=0.6cm, y=0.31cm]
\node[font=\scriptsize, anchor=north east, inner sep=1.5pt] at (axis cs:0,0) {$O$};
\end{axis}
\end{tikzpicture}
\end{minipage}

How did the axis of symmetry help you find a third point?\lines[1]

\wspart{Example 3}{Interpret the Graph of a Quadratic Function}
The graph of the function $f(x)=-10x^2+700x-6{,}000$ shows the profit a company earns for selling headphones at different prices. What is the maximum profit the company can expect to earn?

\begin{minipage}[t]{0.54\textwidth}
\vspace{0pt}
\textbf{Formulate.} The \mbox{$x$-axis} shows selling price and the \mbox{$y$-axis} shows the profit. The \textbf{maximum} $y$-value occurs at the \textbf{vertex} of the parabola.

\textbf{Compute.} $a=-10$, \ $b=700$.\\[2pt]
$h=-\dfrac{b}{2a}=-\dfrac{700}{2(-10)}=-\dfrac{700}{-20}=35$\\[4pt]
$y=-10(35)^2+700(35)-6{,}000$\\
$\phantom{y}=-12{,}250+24{,}500-6{,}000=6{,}250$\\[2pt]
The vertex is $(35,\ 6{,}250)$.

\textbf{Interpret.} A selling price of \$35 per item gives the maximum profit of \$6{,}250.

\small Careful: there is a negative in the formula \emph{and} a negative value of $a$.
\end{minipage}\hfill
\begin{minipage}[t]{0.43\textwidth}
\vspace{0pt}
\begin{tikzpicture}
\begin{axis}[sheetgrid, xmin=0, xmax=90, ymin=0, ymax=8000,
  xtick={0,20,40,60,80}, ytick={0,2000,4000,6000,8000}, minor xtick={0,10,...,90}, minor ytick={0,1000,...,8000},
  x=0.062cm, y=0.00075cm, scaled y ticks=false,
  yticklabels={$0$,$2{,}000$,$4{,}000$,$6{,}000$,$8{,}000$},
  xlabel={Selling Price (\$)}, ylabel={Profit (\$)},
  xlabel style={at={(axis description cs:0.5,-0.1)}, anchor=north, font=\scriptsize},
  ylabel style={at={(axis description cs:-0.24,0.5)}, anchor=south, rotate=90, font=\scriptsize}]
\addplot[very thick, orange!90!black, domain=9.5:60.5, samples=80] {-10*x^2+700*x-6000};
\node[closed] at (axis cs:35,6250) {};
\node[closed] at (axis cs:10,0) {};
\node[closed] at (axis cs:60,0) {};
\node[font=\scriptsize, fill=white, inner sep=1pt, anchor=south] at (axis cs:35,6550) {$(35,\,6{,}250)$};
\end{axis}
\end{tikzpicture}
\end{minipage}

\textbf{Say it:} $a=-10$ is negative, so the parabola opens \blank[1.4cm] and the \textbf{vertex} is a \blank[2cm]. Its $x$-coordinate, \blank[1cm], is the \blank[2.4cm]; its $y$-coordinate, \blank[1.4cm], is the \blank[2.6cm].

\newpage

\wspart{Try It 3}{Maximum height}
A water balloon is tossed straight up from the ground. The height of the water balloon over time can be modeled by the function $y=-14x^2+34x+3$. What was the maximum height of the water balloon after it was tossed?

\begin{minipage}[t]{0.52\textwidth}
\vspace{0pt}
$a=\blank[1cm]$ \quad $b=\blank[1cm]$ \quad $c=\blank[1cm]$

\vspace{6pt}
$h=-\dfrac{b}{2a}=\blank[3.6cm]$
\lines[2]

$y=\blank[5cm]$
\lines[2]

Maximum height: about \blank[2cm] ft, reached after about \blank[1.6cm] s.
\end{minipage}\hfill
\begin{minipage}[t]{0.45\textwidth}
\vspace{0pt}
\begin{tikzpicture}
\begin{axis}[sheetgrid, xmin=0, xmax=4, ymin=0, ymax=25,
  xtick={0,1,2,3}, ytick={0,5,10,15,20}, minor xtick={0,0.5,...,4}, minor ytick={0,2.5,...,25},
  x=1.3cm, y=0.21cm,
  xlabel={Time (seconds)}, ylabel={Height (feet)},
  xlabel style={at={(axis description cs:0.5,-0.1)}, anchor=north, font=\scriptsize},
  ylabel style={at={(axis description cs:-0.12,0.5)}, anchor=south, rotate=90, font=\scriptsize}]
\addplot[very thick, red, domain=0:2.514, samples=80, ->] {-14*x^2+34*x+3};
\end{axis}
\end{tikzpicture}
\end{minipage}

Is the vertex a maximum or a minimum? How do you know before you graph?\lines[2]

% R5 parallel of quiz item 5 (Savvas 2-2 Lesson Quiz item 5 = short-form Q3: graph the diver's path
% y=-7x^2+5x+16 from the board to the water). Same structure, different numbers and setting.
\wspart{Practice}{From the throw to the ground}
A ball is thrown upward from a balcony. Its height, in meters, after $x$ seconds is $y=-5x^2+10x+15$. Graph the function from the point where the ball was thrown to the point where it hits the ground.

\begin{minipage}[t]{0.52\textwidth}
\vspace{0pt}
\textbf{a.} Starting point ($y$-intercept): $(\blank[0.8cm],\blank[0.8cm])$

\vspace{6pt}
\textbf{b.} Vertex: $h=\blank[2.4cm]$, \ $k=\blank[2.4cm]$

\vspace{6pt}
\textbf{c.} Symmetric point: $(\blank[0.8cm],\blank[0.8cm])$

\vspace{6pt}
\textbf{d.} Keep going until the height is 0:

\renewcommand{\arraystretch}{1.5}
\begin{tabular}{|c|c|c|c|c|}
\hline
$x$ & 0.5 & 1.5 & 2.5 & 3 \\ \hline
$y$ & & & & \\ \hline
\end{tabular}

\vspace{6pt}
\textbf{e.} Graph only from the throw to the ground: a closed dot at each end, no arrows.

\textbf{f.} How high is the balcony? What is the greatest height of the ball?\lines[1]
\end{minipage}\hfill
\begin{minipage}[t]{0.45\textwidth}
\vspace{0pt}
\begin{tikzpicture}
\begin{axis}[sheetgrid, xmin=-0.5, xmax=4, ymin=-2, ymax=22,
  xtick={1,2,3,4}, ytick={4,8,12,16,20}, minor xtick={-0.5,0,...,4}, minor ytick={-2,0,...,22},
  x=1.3cm, y=0.24cm,
  xlabel={Time (s)}, ylabel={Height (m)},
  xlabel style={at={(axis description cs:0.5,-0.1)}, anchor=north, font=\scriptsize},
  ylabel style={at={(axis description cs:-0.14,0.5)}, anchor=south, rotate=90, font=\scriptsize}]
\node[font=\scriptsize, anchor=north east, inner sep=1.5pt] at (axis cs:0,0) {$O$};
\end{axis}
\end{tikzpicture}
\end{minipage}

\newpage

\wspart{IXL}{Practice for the quiz}
Go to IXL, Algebra 2, and type the code in the search box. Do the starred skills first.

\vspace{2pt}
{\small\begin{tabularx}{\textwidth}{l Y >{\raggedright\arraybackslash}p{6.4cm}}
\toprule
\textbf{Code} & \textbf{Skill} & \textbf{Feeds} \\
\midrule
\ixlrow{L8C $\bigstar$}{Characteristics of quadratic functions: equations}{Quiz 3 and bonus Quiz 2 (Example 3)}
\ixlrow{S9G $\bigstar$}{Graph a quadratic function}{Quiz 5: the diver's path (Example 2)}
\ixlrow{WMS}{Characteristics of quadratic functions: graphs}{Warm-up: key features off a graph}
\ixlrow{8GN}{Multiply polynomials}{Warm-up: expand $2(x-4)^2+5$}
\bottomrule
\end{tabularx}}

\textbf{SmartScore 70 or above on each starred ($\bigstar$) skill replaces your quiz score.}

% R3: the lesson's DOK 3 tied to a kept Example (Example 3): Practice & Problem Solving 26.
\wspart{Exit question}{Practice 26 (DOK 3)}
A college's business office found the relationship between the number of admissions counselors they employ and the college's profit from tuition could be modeled by the function\newline $y=-10x^2+1{,}500x-35{,}000$.\par

\begin{minipage}[t]{0.55\textwidth}
\vspace{0pt}
\textbf{a.} Graph the function.

\textbf{b.} How many admissions counselors should the college employ to maximize its profit?\lines[1]

\textbf{c.} What is the maximum amount of profit the college can make?\lines[1]

\vspace{4pt}
\textbf{Finish the sentences.}

The \textbf{vertex} is $(\blank[1.2cm],\ \blank[1.8cm])$. The parabola opens \blank[1.4cm] because $a$ is \blank[1.6cm], so the vertex is a \blank[1.8cm].

The college should employ \blank[1.2cm] counselors, the \mbox{$x$-coordinate} of the vertex. The maximum profit is \blank[2cm], the \mbox{$y$-coordinate}.
\end{minipage}\hfill
\begin{minipage}[t]{0.42\textwidth}
\vspace{0pt}
\begin{tikzpicture}
\begin{axis}[sheetgrid, xmin=-50, xmax=200, ymin=-60000, ymax=40000,
  xtick={50,100,150}, ytick={-40000,-20000,20000}, minor xtick={-50,-25,...,200}, minor ytick={-60000,-50000,...,40000},
  x=0.027cm, y=0.000066cm, scaled y ticks=false,
  yticklabels={$-40{,}000$,$-20{,}000$,$20{,}000$}]
\node[font=\scriptsize, anchor=north east, inner sep=1.5pt] at (axis cs:0,0) {$O$};
\end{axis}
\end{tikzpicture}
\end{minipage}

\end{document}
